\documentclass[12pt, letterpaper]{article}

\usepackage[utf8]{inputenc}    
\usepackage[T1]{fontenc}
\usepackage{amsmath, amssymb}  
\usepackage{graphicx}          
\usepackage{hyperref}     
\usepackage{xcolor}     
\usepackage{booktabs}
\usepackage{float}
\usepackage{algorithm}
\usepackage{algpseudocode}
\usepackage{subcaption}
\usepackage{xcolor}
\usepackage[a4paper, margin=2.0cm]{geometry}

\title{Report 10 -- Confronto QALS - GUROBI}

\begin{document}

\maketitle

\newpage

\section{Limiti Gurobi}
Per quanto riguarda i limiti dell'ottimizzatore Gurobi, nella documentazione 
ufficiale {} si può notare come, oltre alle limitazioni imposte dalle possibili licenze,
l'unico reale limite è relativo alla potenza del dispotivo che si sta usando. Ne conviene
che non risulta essere adeguato definire un confronto in relazione a ciò, ma piuttosto
è necessario discutere in relazione al tempo impiegato.
Sulla base dei seguenti risultati dati dalla media di 10 run:
\begin{table}[H]
    \centering
    \begin{tabular}{ccc}
        \toprule
        \textbf{N. items} & \textbf{QALS (s)} & \textbf{Gurobi (s)} \\
        \midrule
        72 & \textcolor{red}{22.42} & 21.83 \\
        73 & \textcolor{red}{21.71} & 21.40 \\
        74 & 19.73 & \textcolor{red}{30.53} \\
        75 & 23.12 & \textcolor{red}{25.12} \\
        76 & 22.47 & \textcolor{red}{77.87} \\
        77 & 24.4 &  \textcolor{red}{88.4} \\
        \bottomrule
    \end{tabular}
    \caption{Confronto dei tempi di esecuzione tra QALS e Gurobi al variare del numero di item.}
    \label{tab:qals_vs_gurobi}
\end{table}
si può concludere che non sempre vale la pena eseguire QALS: come visto nei risultati
fino a 73 items (incluso) ci converrebbe utlizzare Gurobi, dal 74-esimo in poi QALS sembra 
impiegarci meno tempo.

\end{document}